\%============================================================
\chapter{Bài Học: Gieo Điều Tốt, Gặp Điều Tốt (Sow Goodness, Reap Goodness)}
\begin{center}
  \textit{\color{truyengold}``The hand that gives roses always retains a little of the fragrance.''}\\[4pt]
  \textit{(Bàn tay trao tặng hoa hồng, hương thơm vẫn còn lưu lại.)}\\[4pt]
  \small\color{truyenblue}--- Cổ Kim Đông Tây ---
\end{center}
\ngancach

\section{Người Nông Dân Cứu Chàng Trai Chết Đuối}
\begin{truyen}{Người Nông Dân Cứu Chàng Trai Chết Đuối}{Giai thoại về Thủ tướng Churchill --- Anh}
\chuhoa{M}{ột} buổi chiều mùa hè, người nông dân nghèo Hugh Fleming đang làm việc trên cánh đồng thì nghe tiếng kêu cứu thảm thiết vọng ra từ đầm lầy gần đó.\\[4pt]
\textit{(One summer afternoon, poor farmer Hugh Fleming was working in his field when he heard desperate cries for help coming from a nearby swamp.)}

Ông lao xuống ngay không chút do dự và cứu được một cậu bé đang chìm dần trong bùn lầy.\\[4pt]
\textit{(He plunged in without hesitation and saved a young boy who was slowly sinking into the mud.)}

Hôm sau, một thứ gì đó quan trọng hơn đã xảy ra: một chiếc xe ngựa sang trọng đến trước nhà ông, và bên trong là một quý ông giàu có.\\[4pt]
\textit{(The next day, something even more significant occurred: a handsome carriage arrived at his home, carrying a wealthy gentleman.)}

Quý ông đó chính là cha của cậu bé được cứu, ông ngỏ lời muốn cho con trai của người nông dân được học tại trường y khoa danh giá nhất London.\\[4pt]
\textit{(That gentleman was the father of the rescued boy, and he offered to send the farmer's son to study at the most prestigious medical school in London.)}

Người nông dân nhận lời; con trai ông sau này trở thành nhà khoa học Alexander Fleming, người phát minh ra Penicillin cứu sống hàng triệu người.\\[4pt]
\textit{(The farmer accepted; his son later became scientist Alexander Fleming, who invented Penicillin saving millions of lives.)}

Còn cậu bé được cứu khỏi đầm lầy ngày ấy --- chính là người con trai của quý ông kia --- sau này trở thành Thủ tướng Winston Churchill vĩ đại.\\[4pt]
\textit{(And the boy saved from the swamp that day --- the gentleman's own son --- later became the great Prime Minister Winston Churchill.)}

\end{truyen}
\begin{baihoc}
  Một hành động tốt bình thường nhất có thể khởi động một chuỗi nhân quả vĩ đại thay đổi cả lịch sử nhân loại.\\[4pt]
  \textit{(The most ordinary good deed can trigger a chain of karma that changes the entire course of human history.)}
\end{baihoc}

\section{Bát Cháo Của Bà Lão Và Tướng Quân}
\begin{truyen}{Bát Cháo Của Bà Lão Và Tướng Quân}{Truyện Dân Gian Việt Nam}
\chuhoa{T}{rong} một mùa đông giá rét, một bà lão nghèo sống một mình ven rừng đã chia sẻ bát cháo cuối cùng của mình cho một người lính đói khổ lết vào cửa nhà bà.\\[4pt]
\textit{(In a bitter winter, a poor old woman living alone at the forest's edge shared her last bowl of porridge with a starving soldier who dragged himself to her door.)}

Người lính cảm tạ và tiếp tục hành trình gian nan của mình trong cuộc chiến.\\[4pt]
\textit{(The soldier expressed gratitude and continued his arduous wartime journey.)}

Mười năm sau, một đoàn kỵ binh oai nghiêm dừng lại trước căn nhà tranh của bà.\\[4pt]
\textit{(Ten years later, a majestic cavalry unit stopped before the woman's thatched cottage.)}

Vị tướng dẫn đầu xuống ngựa, cúi đầu chào người phụ nữ già nua và trao cho bà một túi vàng.\\[4pt]
\textit{(The general leading the group dismounted, bowed respectfully to the aged woman, and presented her with a bag of gold.)}

Ông nói: 'Thưa cụ, con là người lính đói rét ngày ấy. Bát cháo của cụ đã cho con sức mạnh để tiếp tục sống và chiến đấu. Hôm nay con trở về để báo ơn.'\\[4pt]
\textit{(He said: 'Elder, I am that cold and hungry soldier from years ago. Your bowl of porridge gave me the strength to keep living and fighting. Today I have returned to repay the debt.')}

\end{truyen}
\begin{baihoc}
  Lòng tốt không bao giờ mất đi, nó chỉ đi một vòng dài rồi tìm về đúng người đã gieo.\\[4pt]
  \textit{(Kindness is never lost; it just makes a long journey before returning to the one who sowed it.)}
\end{baihoc}

\section{Người Thương Gia Và Kẻ Trộm Được Tha Thứ}
\begin{truyen}{Người Thương Gia Và Kẻ Trộm Được Tha Thứ}{Charles Dickens phóng tác --- Anh}
\chuhoa{Đ}{êm} đông lạnh giá, một người thương gia tên Edward Montague bắt gặp một cậu bé gầy guộc đang trộm một ổ bánh mì trước tiệm của ông.\\[4pt]
\textit{(On a cold winter night, a merchant named Edward Montague caught a scrawny boy stealing a loaf of bread from his shop.)}

Thay vì giao cậu cho cảnh sát, ông mời cậu vào trong, cho cậu ăn bữa tối ấm áp và hỏi về hoàn cảnh của cậu.\\[4pt]
\textit{(Instead of calling the police, he invited the boy inside, gave him a warm dinner, and asked about his circumstances.)}

Cậu bé kể rằng gia đình đang chết đói vì cha cậu bị bệnh không làm việc được.\\[4pt]
\textit{(The boy explained that his family was starving because his father was too ill to work.)}

Người thương gia tặng cậu thức ăn, thuốc men, và cho cậu một việc làm nhỏ trong cửa tiệm.\\[4pt]
\textit{(The merchant gave the boy food, medicine, and offered him a small job at the shop.)}

Hai mươi năm sau, khi doanh nghiệp của ông sắp sụp đổ vì khủng hoảng, chính người thanh niên đó --- nay là một doanh nhân thành đạt --- đã đến cứu trợ ông.\\[4pt]
\textit{(Twenty years later, when his business was about to collapse due to a crisis, that same young man --- now a successful entrepreneur --- came to his rescue.)}

\end{truyen}
\begin{baihoc}
  Lòng nhân từ đầu tư vào một đứa trẻ tuyệt vọng là khoản đầu tư sinh lợi nhất trên đời.\\[4pt]
  \textit{(Investing kindness in a desperate child is the most profitable investment in the world.)}
\end{baihoc}

\section{Hòn Đá Cản Đường Và Túi Vàng Dưới Lòng Đất}
\begin{truyen}{Hòn Đá Cản Đường Và Túi Vàng Dưới Lòng Đất}{Truyện Ngụ Ngôn Ấn Độ}
\chuhoa{N}{hà} Vua Bharata muốn thử lòng dân chúng, bèn sai người chôn một tảng đá lớn giữa con đường chính, cạnh đó để một túi vàng nhỏ và tờ giấy ghi: 'Ai dời hòn đá này sẽ nhận được túi vàng.'\\[4pt]
\textit{(King Bharata wanted to test his people, so he had a large stone buried in the middle of the main road, with a small bag of gold and a note beside it reading: 'Whoever moves this stone shall receive the bag of gold.')}

Hàng trăm người qua lại, tất cả đều chỉ nhìn hòn đá rồi chê nó quá nặng, rồi bước đi than thở về sự cản trở.\\[4pt]
\textit{(Hundreds of people passed by, and all of them just looked at the stone, complained it was too heavy, then walked away grumbling about the obstruction.)}

Chỉ có một người nông dân già tên Rajan dừng lại, không nói gì, lặng lẽ dùng hết sức bình sinh để đẩy hòn đá sang bên vệ đường.\\[4pt]
\textit{(Only an old farmer named Rajan stopped, said nothing, and silently used all his might to push the stone to the side of the road.)}

Dưới hòn đá, ông tìm thấy túi vàng cùng một tờ giấy thứ hai ghi: 'Người dời hòn đá này không làm vì vàng, mà làm vì ích lợi cho người đi đường sau --- đó là người đáng được thưởng nhất.'\\[4pt]
\textit{(Under the stone, he found the bag of gold and a second note reading: 'The one who moved this stone did not do it for the gold, but for the benefit of those who travel after --- that is the one most deserving of reward.')}

\end{truyen}
\begin{baihoc}
  Hành động vì người khác mà không tính đến lợi ích cá nhân chính là hành động cao thượng nhất và thường được đền đáp xứng đáng nhất.\\[4pt]
  \textit{(Acting for others without calculating personal gain is the noblest act and often the one most fittingly rewarded.)}
\end{baihoc}

\section{Lời Chúc Phúc Của Bà Mẹ Nghèo}
\begin{truyen}{Lời Chúc Phúc Của Bà Mẹ Nghèo}{Truyện dân gian Hàn Quốc}
\chuhoa{N}{gày} xưa ở Hàn Quốc, có một bà mẹ nghèo mù lòa nuôi hai đứa con mà không có tiền mua đồ ăn trong mùa đông.\\[4pt]
\textit{(Long ago in Korea, there was a poor blind mother raising two children with no money to buy food in winter.)}

Một thương gia qua đường thấy cảnh tượng đó, lặng lẽ để lại một bao gạo trước cửa nhà bà rồi đi.\\[4pt]
\textit{(A passing merchant saw the scene, quietly left a sack of rice at her door, and departed.)}

Bà mẹ không biết ai đã làm điều đó, bà chỉ biết ngước mặt lên trời cầu nguyện, ban phúc cho người đã giúp bà.\\[4pt]
\textit{(The mother did not know who had done it; she only looked up to heaven and prayed, blessing whoever had helped her.)}

Người thương gia đó, tên là Park Soo-Il, suốt cuộc đời làm ăn luôn gặp may mắn kỳ lạ; mỗi khi gặp nguy hiểm đều thoát nạn một cách khó hiểu.\\[4pt]
\textit{(That merchant, named Park Soo-Il, throughout his business life encountered strange good fortune; whenever he faced danger, he escaped in inexplicable ways.)}

Ông nghĩ đó là lời cầu nguyện của người mẹ mù nghèo ngày ấy vẫn đang che chở cho ông.\\[4pt]
\textit{(He believed those were the prayers of that poor blind mother from long ago still shielding him.)}

\end{truyen}
\begin{baihoc}
  Lời cầu chúc chân thành từ trái tim người được cứu giúp là lá bùa hộ mạng mạnh nhất trên thế gian.\\[4pt]
  \textit{(A sincere blessing from the heart of someone helped is the most powerful protective talisman in the world.)}
\end{baihoc}

\section{Người Bác Sĩ Và Bệnh Nhân Không Có Tiền}
\begin{truyen}{Người Bác Sĩ Và Bệnh Nhân Không Có Tiền}{Albert Schweitzer --- Gabon, châu Phi}
\chuhoa{Ở}{} vùng Lambaréné thuộc Gabon Trung Phi, bác sĩ Albert Schweitzer xây bệnh viện và chữa bệnh miễn phí cho người dân nghèo trong suốt hơn 50 năm.\\[4pt]
\textit{(In the Lambaréné region of Central Africa's Gabon, Doctor Albert Schweitzer built a hospital and treated the poor for free for over 50 years.)}

Ông đã từ bỏ sự nghiệp âm nhạc và triết học lừng danh ở châu Âu để làm điều này.\\[4pt]
\textit{(He gave up a celebrated career in music and philosophy in Europe to do this.)}

Nhiều người hỏi ông: 'Tại sao ông lại hi sinh như vậy?' Ông chỉ trả lời bằng nụ cười: 'Bởi vì tôi may mắn hơn họ.'\\[4pt]
\textit{(Many asked him: 'Why do you sacrifice like this?' He only answered with a smile: 'Because I am more fortunate than they are.')}

Năm 1952, ông được trao giải Nobel Hòa bình; ông dùng toàn bộ số tiền thưởng để mở rộng bệnh viện.\\[4pt]
\textit{(In 1952, he was awarded the Nobel Peace Prize; he used the entire prizemoneu to expand the hospital.)}

Ông viết trong nhật ký: 'Tôi không cho đi vì muốn nhận lại. Nhưng cuộc đời lạ lùng thay --- càng cho đi tôi càng nhận được nhiều hơn, không phải tiền bạc, mà là sự bình yên trong tâm hồn.'\\[4pt]
\textit{(He wrote in his diary: 'I do not give in order to receive. But how strange life is --- the more I give, the more I receive, not money, but peace within the soul.')}

\end{truyen}
\begin{baihoc}
  Người sống vì người khác không phải người mất mát nhiều nhất --- họ là người giàu có nhất vì sở hữu thứ mà tiền bạc không mua nổi.\\[4pt]
  \textit{(Those who live for others are not those who lose the most --- they are the richest because they own what money cannot buy.)}
\end{baihoc}

\section{Cây Đại Thụ Và Người Du Khách Mệt Mỏi}
\begin{truyen}{Cây Đại Thụ Và Người Du Khách Mệt Mỏi}{Truyện ngụ ngôn Hy Lạp cổ}
\chuhoa{H}{ai} người lữ hành kiệt sức đang đi trên con đường dài dưới trời nắng chói chang thì nhìn thấy một cây đại thụ cao lớn tỏa bóng mát rộng.\\[4pt]
\textit{(Two exhausted travelers walking a long road under a blazing sun spotted a tall great tree casting wide shade.)}

Họ ngồi xuống nghỉ ngơi bên dưới, hít thở không khí mát lành và cảm thấy sức lực hồi phục.\\[4pt]
\textit{(They sat down to rest beneath it, breathed the cool air, and felt their strength restore.)}

Một người đàn ông nói: 'Cái cây này vô dụng quá, không có trái, không có gỗ tốt, chỉ có bóng mát mà thôi.'\\[4pt]
\textit{(One man said: 'This tree is so useless, no fruit, no valuable timber, only shade.')}

Người kia lại nói: 'Anh nói sai rồi. Chính cái bóng mát mà anh cho là vô dụng đó đang cứu chúng ta khỏi chết nắng lúc này.'\\[4pt]
\textit{(The other replied: 'You are wrong. It is precisely this shade you call useless that is saving us from heatstroke right now.')}

Cây đại thụ lúc đó mới rung lá thì thầm: 'Điều tốt lành mà anh không nhìn thấy giá trị ngay lúc này, có thể đang cứu anh mà anh không hay.'\\[4pt]
\textit{(The great tree then rustled its leaves and whispered: 'The goodness whose value you cannot see at this moment may be saving you without your knowledge.')}

\end{truyen}
\begin{baihoc}
  Không phải mọi ân huệ đều xuất hiện dưới dạng ta nhận ra ngay; đôi khi điều tốt lành đang bao quanh ta mà ta chẳng hề hay biết.\\[4pt]
  \textit{(Not all blessings appear in forms we instantly recognize; sometimes goodness surrounds us without our knowledge.)}
\end{baihoc}

\section{Người Chèo Đò Và Triết Gia}
\begin{truyen}{Người Chèo Đò Và Triết Gia}{Truyện Phật Giáo --- Ấn Độ}
\chuhoa{M}{ột} triết gia nổi tiếng hỏi người chèo đò già: 'Ông biết thiên văn không?' Người chèo đò lắc đầu. 'Triết học không?' Lắc đầu. 'Văn học không?' Vẫn lắc đầu.\\[4pt]
\textit{(A famous philosopher asked the old ferryman: 'Do you know astronomy?' The ferryman shook his head. 'Philosophy?' Shook his head. 'Literature?' Still shook his head.)}

Nhà triết học chê bai: 'Ôi, ông đã mất đi phần lớn cuộc đời rồi!'\\[4pt]
\textit{(The philosopher mocked: 'Oh, you have wasted most of your life!')}

Giữa đường, con thuyền bị sóng đánh và bắt đầu chìm. Người chèo đò hỏi: 'Ông có biết bơi không?'\\[4pt]
\textit{(In the middle of the river, waves struck the boat and it began to sink. The ferryman asked: 'Can you swim?')}

Nhà triết học lắc đầu hoảng loạn. Người chèo đò bình thản nói: 'Thế thì ông mới thực sự mất cả cuộc đời.'\\[4pt]
\textit{(The philosopher shook his head in panic. The ferryman calmly said: 'Then it is you who has truly wasted your whole life.')}

Khi người chèo đò vớt nhà triết học lên bờ an toàn, ông nói thêm: 'Mỗi người có loại kiến thức của riêng mình. Hãy tôn trọng điều đó trước khi phán xét.'\\[4pt]
\textit{(When the ferryman pulled the philosopher safely to shore, he added: 'Each person has their own kind of knowledge. Respect that before judging.')}

\end{truyen}
\begin{baihoc}
  Sự khiêm tốn là tiền lệ của sự học hỏi; kẻ coi thường người khác đang đóng cánh cửa với những bài học quý giá nhất.\\[4pt]
  \textit{(Humility is the prerequisite of learning; those who look down on others are shutting the door to the most precious lessons.)}
\end{baihoc}

\section{Người Thợ Rèn Và Viên Gạch Cuối Cùng}
\begin{truyen}{Người Thợ Rèn Và Viên Gạch Cuối Cùng}{Truyện cổ tích Đức}
\chuhoa{M}{ột} người thợ rèn già tên Hans sắp nghỉ hưu sau 40 năm làm việc cần cù và trung thực.\\[4pt]
\textit{(An old blacksmith named Hans was about to retire after 40 years of diligent and honest work.)}

Ngày cuối cùng trước khi rời xưởng, ông tìm thấy một đứa trẻ mồ côi đang co ro dưới mưa trước cổng xưởng.\\[4pt]
\textit{(On his last day before leaving the workshop, he found an orphaned child huddled in the rain at the workshop gate.)}

Dù đã mệt mỏi, ông vẫn dẫn cậu vào trong, cho cậu ăn, và dùng những đồng tiền tiết kiệm cuối cùng để mua cho cậu một bộ quần áo ấm.\\[4pt]
\textit{(Despite his exhaustion, he brought the boy inside, fed him, and used his last saved coins to buy the boy warm clothes.)}

Mọi người trong làng bảo ông: 'Hans, ông điên rồi, đó là tiền dưỡng già của ông.'\\[4pt]
\textit{(Villagers told him: 'Hans, you're mad, that was your retirement money.')}

Hans chỉ mỉm cười: 'Tôi đã xây ngôi nhà này bằng 40 năm lao động, nhưng viên gạch quan trọng nhất là viên gạch cuối cùng tôi đặt xuống hôm nay.'\\[4pt]
\textit{(Hans only smiled: 'I built this house with 40 years of labor, but the most important brick is the one I lay down today.')}

Ba năm sau, cậu bé đã trở thành học trò của một thương gia và tìm đến trả ơn người thợ rèn già bằng cách chăm sóc ông suốt những năm tuổi già.\\[4pt]
\textit{(Three years later, the boy had become a merchant's apprentice and came back to repay the old blacksmith by caring for him through all his elderly years.)}

\end{truyen}
\begin{baihoc}
  Hành động tốt ở phút cuối của một hành trình có thể là hành động vĩ đại nhất và đơm hoa kết quả ngọt ngào nhất.\\[4pt]
  \textit{(A good deed at the final moment of a journey may be the greatest act and the one bearing the sweetest fruit.)}
\end{baihoc}

\section{Bãi Sao Biển Và Cô Bé}
\begin{truyen}{Bãi Sao Biển Và Cô Bé}{Truyện ngụ ngôn hiện đại --- Loren Eiseley}
\chuhoa{S}{áng} sớm trên bãi biển sau cơn bão, hàng nghìn con sao biển bị sóng ném lên bờ đang chết dần trong ánh nắng.\\[4pt]
\textit{(Early morning on the beach after a storm, thousands of starfish thrown ashore by the waves were slowly dying in the sun.)}

Một cô bé đang cúi xuống nhặt từng con sao biển một rồi ném chúng trở lại đại dương.\\[4pt]
\textit{(A little girl was bending down picking up each starfish one by one and throwing them back into the ocean.)}

Một ông lão đi qua hỏi: 'Cháu làm gì vậy? Có hàng nghìn con sao biển, cháu không thể cứu tất cả được đâu.'\\[4pt]
\textit{(An old man passing by asked: 'What are you doing? There are thousands of starfish, you cannot save them all.')}

Cô bé nhặt thêm một con, ném xuống biển và nói: 'Nhưng con cứu được con này.'\\[4pt]
\textit{(The girl picked up another one, threw it into the sea and said: 'But I saved this one.')}

Ông lão đứng lặng một lúc, rồi cúi xuống nhặt một con sao biển khác và ném ra biển.\\[4pt]
\textit{(The old man stood quietly for a moment, then bent down, picked up another starfish and tossed it into the sea.)}

Hai người cứ thế tiếp tục, và cùng nhau họ cứu được hơn một ngàn con sao biển trong buổi sáng hôm đó.\\[4pt]
\textit{(The two continued like this, and together they saved more than a thousand starfish that morning.)}

\end{truyen}
\begin{baihoc}
  Đừng để cái chưa làm được ngăn cản bạn làm điều bạn đang làm được; mỗi hành động tốt dù nhỏ bé đều có ý nghĩa với ai đó.\\[4pt]
  \textit{(Do not let what you cannot do prevent you from doing what you can; every good act, however small, matters to someone.)}
\end{baihoc}
